\chapter{Step 09: The ULM on ICE}
\label{ch:step09}

% A register of the ULM, e.g. \reg{1} gives %1.
\newcommand{\reg}[1]{\texttt{\%#1}}

This is the last step: the ULM itself. Almost all of its parts are already
there. Chapter~\ref{ch:step05} built a small ALU with a register file and a
decoder, Chapter~\ref{ch:step07} the serial I/O with its FIFOs,
Chapter~\ref{ch:step08} the RAM and the loader that fills it with a program
typed into a terminal. What is new is the part that makes these parts a
computer: a \term{control unit} that fetches an instruction from the RAM,
decodes it, and lets the ALU, the register file, the RAM or the I/O execute
it, over and over again. Everything is in \cc{09_ulm_on_ice}; the files that
are links to earlier steps are unchanged.

\section{The Instruction Set}
\label{sec:step09-isa}

Every ULM instruction is a 32-bit word. Its top byte, bits 31 to 24, is the
\term{opcode}; it says what to do. The other 24 bits are split into fields,
and how they are split depends on the instruction. The file
\cc{examples/isa.txt} lists these splits as \term{formats}:

\begin{center}
\small
\begin{tabular}{lcccc@{\quad}l}
\toprule
format & 31--24 & 23--20 & 19--16 & 15--0 & used by \\
\midrule
\texttt{U8}      & OP & \multicolumn{2}{c}{\textit{imm} (8 bits)} & & \texttt{halt 0x42}, \texttt{putc 'a'} \\
\texttt{R}       & OP & $x$ & & & \texttt{putc \%x}, \texttt{getc \%x} \\
\texttt{RR}      & OP & $y$ & $x$ & & \texttt{jmp \%y, \%x} \\
\texttt{REL\_JMP} & OP & \multicolumn{3}{c}{\textit{offset} (24 bits, signed)} & \texttt{jz}, \texttt{jmp}, \dots \\
\texttt{U20\_R}  & OP & \textit{dest} & \multicolumn{2}{c}{\textit{imm} (20 bits)} & \texttt{loadz imm, \%dest} \\
\texttt{R\_R\_R} & OP & $z$ & $y$ & $x$ (15--12) & \texttt{addq \%x, \%y, \%z} \\
\texttt{U16\_R\_R} & OP & $z$ & $y$ & \textit{imm} & \texttt{addq imm, \%y, \%z} \\
\texttt{MR\_R}, \texttt{R\_MR} & OP & \textit{data} & \textit{addr} & \textit{offset} (signed) & \texttt{movq 8(\%addr), \%data} \\
\bottomrule
\end{tabular}
\end{center}

Register numbers are 4 bits wide, for the 16 registers \reg{0} to \reg{15}.
\reg{0} is special: it always reads as 0, and writing it has no effect. The
instructions that the hardware understands are these (``imm'' is a constant
from the instruction, $u$ and $s$ mean that it is read as an unsigned or a
signed number):

\begin{center}
\small
\begin{tabular}{@{}l@{\quad}l@{\quad}l@{}}
\toprule
opcode & instruction & effect \\
\midrule
\texttt{01}, \texttt{02} & \texttt{halt imm}, \texttt{halt \%x} & stop, with an exit code \\
\texttt{05} & \texttt{jmp offset} & $\mathit{ip} \gets \mathit{ip} + 4\cdot\mathit{offset}$ \\
\texttt{03}, \texttt{04}, \texttt{06} & \texttt{jnz}, \texttt{jz}, \texttt{jb} & the same, if ZF $=0$, ZF $=1$, CF $=1$ \\
\texttt{08}--\texttt{0D} & \texttt{jbe}, \texttt{ja}, \texttt{jae}, \texttt{jge}, \texttt{jnge}, \texttt{jng} & conditional jumps on other flag combinations\\
\texttt{07} & \texttt{jmp \%y, \%x} & $\%x \gets \mathit{ip}+4$, $\mathit{ip} \gets \%y$ (\texttt{call}, \texttt{ret}) \\
\texttt{10}, \texttt{11} & \texttt{loadz imm, \%dest}, \texttt{loads} & $\%\mathit{dest} \gets u(\mathit{imm})$, resp.\ $s(\mathit{imm})$ \\
\texttt{12}, \texttt{13} & \texttt{addq imm, \%y, \%z}, \texttt{addq \%x, \%y, \%z} & $\%z \gets \%y + \mathit{imm}$, resp.\ $\%y + \%x$ \\
\texttt{14}, \texttt{15} & \texttt{subq} (same two forms) & $\%z \gets \%y - \mathit{imm}$, resp.\ $\%y - \%x$ \\
\texttt{16}, \texttt{17} & \texttt{mulw} (same two forms) & product of the lowest 16 bits \\
\texttt{18}--\texttt{1D} & \texttt{andq}, \texttt{shrq}, \texttt{shlq} (both forms) & bitwise and, shift right, shift left \\
\texttt{20}, \texttt{21} & \texttt{movzbq offset(\%addr), \%data}, \texttt{movsbq} & load one byte, zero- resp.\ sign-extended \\
\texttt{22} & \texttt{movb \%data, offset(\%addr)} & store one byte \\
\texttt{23}, \texttt{24} & \texttt{movq} (load and store) & the same for 8 bytes \\
\texttt{30}, \texttt{31} & \texttt{putc \%x}, \texttt{putc imm} & send a character \\
\texttt{32} & \texttt{getc \%x} & wait for a character, $\%x \gets$ its code \\
\bottomrule
\end{tabular}
\end{center}

All arithmetic is modulo $2^{64}$. The flags are set by the last ALU
instruction: ZF if its result was 0, SF to its bit 63, CF if an addition had a
carry or a subtraction a borrow, OF if it overflowed as a signed operation (the
other ALU instructions clear CF and OF), so \texttt{subq \%x, \%y, \%0} followed by \texttt{jb}
jumps if $u(\%y) < u(\%x)$. Since a jump offset counts instructions, not
bytes, the hardware multiplies it by 4.

The file \cc{examples/isa.txt} describes exactly these instructions, in the
input format of the \cc{ulm-generator}: first the formats, then for each
opcode its format, its assembler notation, and its effect as a few lines of C.
From this file the \cc{ulm-generator} builds an assembler (\cc{ulmas}), which
translates the \cc{.s} files of \cc{examples} into machine code (the files
without extension contain the same bytes as hex digits), and a simulator of
the ULM in software (\cc{ulm}), in which a program can be tried out before it
goes to the board.

\section{The Architecture}
\label{sec:step09-arch}

Figure~\ref{fig:step09-blocks} shows how \cc{test.sv} connects the parts. The
I/O device \cc{dev_io} and the RAM \cc{dev_ram} have two possible users: the
loader of Chapter~\ref{ch:step08} and the new control unit \cc{dev_cu}. The
switches \cc{dev_io_switch} and \cc{dev_ram_switch} are multiplexers that
connect either the one or the other, and they are controlled by the loader's
output \cc{done}:

\svpart{firstline=98,lastline=123}{09_ulm_on_ice/test.sv}

The same signal enables the control unit (\cc{.en(loader_done)}). So the
board is in one of two modes. As long as \cc{loader_done} is 0, the loader
owns I/O and RAM, reads hex digits and writes the bytes to the RAM, starting
at address 0; the control unit stands still. When the loader receives the
character with code 4 (end of transmission, Ctrl-D), it sets \cc{done}, the
switches flip, and the control unit starts to execute the program at
address 0. The button \cc{BTN1} is the reset: it clears \cc{done}, so the
loader takes over again, and puts the control unit back to its start state.

\begin{figure}[ht]
\centering
\begin{tikzpicture}[x=0.93cm,y=0.93cm,
    blk/.style={box,fill=blue!7,minimum height=12mm,minimum width=20mm},
    sw/.style={box,fill=black!5,minimum width=10mm,minimum height=32mm}]
  \node[blk] (io) at (0,0) {\cc{dev_io}\\[2pt]\footnotesize UART,\\\footnotesize 2 FIFOs};
  \node[sw] (iosw) at (2.5,0) {\rotatebox{90}{\cc{dev_io_switch}}};
  \node[blk] (ld) at (6.1,1.6) {\cc{dev_loader}};
  \node[box,fill=green!8,minimum width=38mm,minimum height=24mm] (cu) at (6.1,-1.6) {};
  \node[font=\small,anchor=north] at (cu.north) {\cc{dev_cu}};
  \node[box,fill=white,font=\footnotesize] at (5.0,-1.85) {state\\machine,\\\cc{cu_ip}};
  \node[box,fill=white,font=\footnotesize] at (7.25,-1.85) {\cc{dev_decoder}};
  \node[sw] (ramsw) at (9.7,0) {\rotatebox{90}{\cc{dev_ram_switch}}};
  \node[blk] (ram) at (12.1,0) {\cc{dev_ram}\\[2pt]\footnotesize 4 SPRAMs,\\\footnotesize 128\,KiB};
  \node[blk] (alu) at (4.9,-4.4) {\cc{dev_alu}\\[1pt]\footnotesize + DSP};
  \node[blk] (rf) at (7.4,-4.4) {\cc{dev_reg_file}\\[1pt]\footnotesize 8 block RAMs};
  % pins
  \draw[arrow] (-1.9,0.3) node[left,font=\small] {\cc{RX}} -- (io.west |- 0,0.3);
  \draw[arrow] (io.west |- 0,-0.3) -- (-1.9,-0.3) node[left,font=\small] {\cc{TX}};
  % connections
  \draw[biarrow] (io) -- (iosw);
  \draw[biarrow] (iosw.east |- ld) -- node[above,note] {\cc{io_loader}} (ld);
  \draw[biarrow] (iosw.east |- 0,-1.2) -- node[above,note] {\cc{io_cpu}} (cu.west |- 0,-1.2);
  \draw[biarrow] (ld) -- node[above,note] {\cc{ram_loader}} (ramsw.west |- ld);
  \draw[biarrow] (cu.east |- 0,-1.2) -- node[above,note] {\cc{ram_cpu}} (ramsw.west |- 0,-1.2);
  \draw[biarrow] (ramsw) -- (ram);
  \draw[biarrow] (alu.north) -- (alu.north |- cu.south);
  \draw[biarrow] (rf.north) -- (rf.north |- cu.south);
  % select
  \draw[arrow,dashed,draw=red!60!black] (ld.north) -- ++(0,0.5) -| node[pos=0.25,above,note] {\cc{loader_done}: select, and \cc{en} of \cc{dev_cu}} (ramsw.north);
  \draw[arrow,dashed,draw=red!60!black] (ld.north) ++(0,0.5) -| (iosw.north);
  % display
  \node[note,anchor=north,text width=36mm] at (11.4,-2.4)
    {hex display: the byte being loaded, later the exit code. LEDs: input buffer empty/full, loader done, CPU halted.};
\end{tikzpicture}
\caption{The ULM on the FPGA. The loader and the control unit share the I/O
and the RAM through two switches; \cc{loader_done} decides who is
connected. The ALU and the register file are separate modules, used only by
the control unit.}
\label{fig:step09-blocks}
\end{figure}

The arrows in Figure~\ref{fig:step09-blocks} are interfaces with modports, as
in Chapter~\ref{ch:step07}: \cc{if_io}, \cc{if_ram}, and two new ones for the
ALU and the register file. Look at \cc{if_reg_file.sv}: the register file has
one write port (\cc{addr_in}, \cc{data_in}, and \cc{op} to say whether to
write) and two read ports (\cc{addr_out0}, \cc{data_out0} and
\cc{addr_out1}, \cc{data_out1}), because an instruction like \texttt{addq
\%x, \%y, \%z} needs two registers at the same time. The ALU's interface
\cc{if_alu.sv} has the operation \cc{op}, the operands \cc{a} and \cc{b}, the
result \cc{s} and the four flags. The packages \cc{pkg_alu}, \cc{pkg_bus},
\cc{pkg_cu}, \cc{pkg_io} and \cc{pkg_reg} contain the \cc{enum} types for the
operations of each unit, and \cc{pkg_reg} also the sizes:

\svpart{firstline=4,lastline=17}{09_ulm_on_ice/pkg_reg.sv}

This is the trick of \cc{pkg_ram} from Chapter~\ref{ch:step07} again: an
\cc{enum} with chosen values, used as a set of named constants.

\section{The Register File and the ALU}
\label{sec:step09-rf-alu}

The register file is an array of sixteen 64-bit words with synchronous reads:

\svfile{09_ulm_on_ice/dev_reg_file.sv}

As in Chapter~\ref{ch:step07}, a memory array that is read inside a clocked
block becomes block RAM. A block RAM of the iCE40 (\cc{SB_RAM40_4K}) has one
read port and one write port and is at most 16 bits wide. So yosys needs four
of them side by side for 64 bits, and, since we have two read ports, a second
set of four that holds a copy of the same data: both copies are written
together, each is read through its own port. That makes eight of the 30 block
RAMs, of which we use only 16 of 256 entries. \cc{mem[0]} starts as 0, and
writes to address 0 are suppressed, so address 0 always holds 0: that is how
\reg{0} is always 0. The attribute \cc{(* no_rw_check *)} is explained in
Section~\ref{sec:step09-resources}.

The ALU computes all operations in combinational logic, with 64 bits, and
selects one of them by \cc{op}. Unlike the Mini-ALU of
Chapter~\ref{ch:step05}, which uses one adder for addition and subtraction (for a
subtraction it inverts the second operand with $\xor\,\mathit{sub}$ and adds
the carry-in $\mathit{sub}$), it has a separate adder and subtractor:

\svfile{09_ulm_on_ice/dev_alu.sv}

New is the register \cc{acc}. When \cc{op} is not \cc{ALU_NOP}, it takes over
the result at the clock edge, and the flags \cc{cf_acc}, \cc{of_acc} and
\cc{sf_acc} take over theirs. When \cc{op} is \cc{ALU_NOP}, the output
\cc{alu.s} is \cc{acc}. So the ALU remembers its last result, and the zero
flag \cc{alu.zf}, which is computed from \cc{alu.s}, belongs to it. This is
exactly what the conditional jumps need: they look at the flags of the last
ALU instruction. The product for \cc{mulw} uses only the lowest 16 bits of
each operand. yosys maps it to one of the eight \term{DSP blocks} of the
iCE40UP5K (\cc{SB_MAC16}), a ready-made $16\times16$-bit multiplier, so the
multiplication costs no logic cells at all. The netlist shows it as
\cc{dev_alu0.prod_SB_MAC16}.

\section{The Decoder}
\label{sec:step09-decoder}

The decoder in \cc{dev_decoder.sv} takes the instruction word \cc{ir} and
splits it into four groups of signals, one for each unit that may be
involved: \cc{instr_alu}, \cc{instr_bus} (RAM access), \cc{instr_cu} (halt and
jumps), and \cc{instr_io}. Each group is an interface; this is the one for the
ALU:

\svpart{firstline=9,lastline=24}{09_ulm_on_ice/if_instr_alu.sv}

Here is how the decoder fills it:

\svpart{firstline=144,lastline=194}{09_ulm_on_ice/dev_decoder.sv}

Read it as a circuit. The \cc{always_comb} block is the actual decoder: for
the opcode \cc{op = ir[31:24]}, a \cc{case} computes the wires of
\cc{instr_alu_next}, with defaults for all opcodes that are not ALU
instructions (\cc{ALU_NOP}). The register numbers are simply cut out of the
fixed positions of the formats, \cc{ir[23:20]} for $z$, \cc{ir[19:16]} for $y$,
\cc{ir[15:12]} for $x$, whether the instruction uses them or not. The
\cc{always_ff} block is a register for the decoded instruction: when \cc{en}
is 1, it takes over the wires of \cc{instr_alu_next}. The other three groups
are built in the same way. For \cc{loadz} the decoder plays a small trick: it
turns the instruction into an addition $\%\mathit{dest} \gets
\mathit{imm} + \%0$, so the ALU needs no separate ``load'' operation.

Note how the interfaces are used here. \cc{if_instr_alu instr_alu_next();}
creates an interface inside the module, not as a port, only to have a name for
a bundle of wires: \cc{instr_alu_next.op}, \cc{instr_alu_next.s_reg}, \dots{}
In other SystemVerilog code you will often see a \cc{struct} for this purpose:
\begin{lstlisting}
typedef struct packed {
    pkg_alu::op_t op;
    logic [3:0]   s_reg, b_reg, a_reg;
    ...
} instr_alu_t;
\end{lstlisting}
A packed struct is a vector whose bits have names, so a single assignment can
copy all fields at once (\cc{instr_alu <= instr_alu_next;}). The ULM does not
use structs; the interfaces do the same job, at the price of copying the
fields one by one.

The conditional jumps are decided in the decoder too: for \texttt{jz}, the
decoder sets \cc{instr_cu_next.op} to \cc{CU_REL_JMP} if \cc{alu.zf} is 1 and
to \cc{CU_NOP} otherwise. This is why the decoder is connected to the ALU.

\section{The Control Unit}
\label{sec:step09-cu}

The control unit \cc{dev_cu} is a state machine (Chapter~\ref{ch:step06}).
Its states are in \cc{pkg_cu}:

\svpart{firstline=14,lastline=21}{09_ulm_on_ice/pkg_cu.sv}

and its next-state logic is short:

\svpart{firstline=38,lastline=88}{09_ulm_on_ice/dev_cu.sv}

Every instruction takes the same five states, FETCH, DECODE,
LOAD\_OPERANDS, EXECUTE, INCREMENT, one clock cycle each
(Figure~\ref{fig:step09-states}). There are three exceptions. In EXECUTE, a
\texttt{getc} waits until the input buffer has a character
(\cc{io.getc_en}), and a \texttt{putc} waits until \cc{dev_io} has taken
the character (\cc{io.putc_push_done}, see below). And a \texttt{halt} is
noticed only in the FETCH state after it, when the decoded \cc{instr_cu.op}
is a halt; then the machine goes to HALTED and stays there until reset. As long as \cc{en} (that is,
\cc{loader_done}) is 0, the state register keeps its value.

\begin{figure}[ht]
\centering
\begin{tikzpicture}[x=0.93cm,y=0.93cm,
    st/.style={draw,rounded corners=6pt,fill=green!8,minimum height=10mm,minimum width=20mm,font=\small,align=center,inner sep=3pt}]
  \node[st] (f) at (0,0) {FETCH};
  \node[st] (d) at (3.0,0) {DECODE};
  \node[st] (l) at (6.0,0) {LOAD\_\\OPERANDS};
  \node[st] (e) at (9.0,0) {EXECUTE};
  \node[st] (i) at (12.0,0) {INCREMENT};
  \node[st,fill=red!8] (h) at (0,-3.0) {HALTED};
  \draw[arrow] (f) -- (d);
  \draw[arrow] (d) -- (l);
  \draw[arrow] (l) -- (e);
  \draw[arrow] (e) -- (i);
  \draw[arrow] (i.north) -- ++(0,2.0) -| node[pos=0.25,above,note] {always} (f.north);
  \draw[arrow] (f.west) -- ++(-0.4,0) |- node[pos=0.25,left,note] {after\\a halt} (h.west);
  \draw[arrow] ([xshift=-4mm]e.north) .. controls +(-0.3,0.9) and +(0.3,0.9) .. node[above,note,align=center] {\texttt{getc}: no character,\\\texttt{putc}: not taken} ([xshift=4mm]e.north);
  \foreach \x/\t in {0/{RAM address $\gets$ \cc{cu_ip}},
                     3.0/{\cc{cu_ir} from RAM; decoder registers load},
                     6.0/{register file reads operands},
                     9.0/{ALU computes, RAM access starts, or \texttt{putc} sends},
                     12.0/{write result register (if any), update \cc{cu_ip}, \texttt{getc} pops}}
    \node[note,text width=25mm,anchor=north] at (\x,-0.8) {\t};
\end{tikzpicture}
\caption{The state machine of the control unit, with what happens in each
state. Each state lasts one clock cycle, except EXECUTE while \texttt{getc} or
\texttt{putc} waits.}
\label{fig:step09-states}
\end{figure}

Why five cycles? Because every memory in the ULM is synchronous: the SPRAM
and the block RAMs take an address at one clock edge and deliver the data
after it. So the instruction, fetched with the address \cc{cu_ip} in FETCH, is
only available in DECODE; the decoded register numbers, latched at the end of
DECODE, can only then be sent to the register file, whose data appear after
LOAD\_OPERANDS; the ALU's result is caught in \cc{acc} at the end of EXECUTE,
and written back in INCREMENT.

The rest of \cc{dev_cu.sv} is combinational logic that says, for each state
and each kind of instruction, what goes to the RAM, the register file, the
ALU and the I/O. Instructions and data are in the same RAM. A second
\cc{dev_ram_switch}, inside the control unit, gives the RAM to the instruction
fetch in FETCH and DECODE, and to the data access in all other states:

\svpart{firstline=93,lastline=122}{09_ulm_on_ice/dev_cu.sv}

So the very switch that separates loader and CPU is used once more, a nice
example of a module as a reusable circuit. The instruction is fetched as a
\cc{RAM_LONG}, four bytes.

The register file is controlled like this:

\svpart{firstline=216,lastline=251}{09_ulm_on_ice/dev_cu.sv}

This is a multiplexer for the addresses and the write data of the register
file. Its defaults are those of an ALU instruction: read registers \cc{a_reg}
and \cc{b_reg}, write the ALU output \cc{alu.s} to \cc{s_reg}. The
\cc{if}-chain after it replaces them for the instructions that need something
else (the listing shows two of them; further down follow \texttt{halt \%x},
\texttt{jmp \%y, \%x}, and the loads and stores). The register file writes in
the INCREMENT state, and only for an instruction that has a result register:
\cc{writes_reg} is 1 for the ALU instructions, \texttt{getc}, \texttt{jmp \%y,
\%x} (which saves the return address in $\%x$) and the loads. This condition
matters. The register numbers are cut out of the instruction word whether the
instruction uses them or not, so for \texttt{putc 'A'} (\texttt{0x31410000})
\cc{s_reg} is 4, and for \texttt{jmp fetch} (\texttt{0x05FFFFFB}) it is 15.
Without \cc{writes_reg}, these instructions would write the ALU's last result
into \reg{4} and \reg{15} (Exercise~\ref{ex:step09-putc}).

The ALU gets its operands from the register file or, for an immediate
operand, from the decoded instruction, and it operates only in EXECUTE; in
all other states it gets \cc{ALU_NOP} and keeps \cc{acc}:

\svpart{firstline=288,lastline=303}{09_ulm_on_ice/dev_cu.sv}

The I/O instructions are controlled by this block:

\svpart{firstline=187,lastline=214}{09_ulm_on_ice/dev_cu.sv}

\texttt{getc} takes its character from the input buffer in INCREMENT, after
EXECUTE has waited for \cc{io.getc_en}. \texttt{putc} offers its character
in EXECUTE (\cc{io.putc_push}), and the state machine stays in EXECUTE until
\cc{dev_io} answers with \cc{io.putc_push_done} (Chapter~\ref{ch:step07}).
\cc{dev_io} does not always take the character at once: not while the output
buffer is full, and not in a cycle in which it pushes the echo of a received
byte. In that case \cc{putc_push_done} stays 0, and the control unit simply
offers the character again in the next cycle. When \cc{dev_io} takes it, it
sets \cc{putc_push_done} for one cycle; in that cycle \cc{putc_push} is still
1, but \cc{dev_io} is busy with its own push and ignores it, so each character
goes into the buffer exactly once. Without this handshake, a program that
sends more than the 512 bytes the output buffer holds would lose characters.

The instruction pointer \cc{cu_ip} is a 17-bit register (the RAM has
$2^{17}$ bytes). In INCREMENT it becomes $\mathit{ip}+4$, or $\mathit{ip} +
4\cdot\mathit{offset}$ for a jump that is taken (the decoder already appended
two zero bits to the offset), or the content of a register for \texttt{jmp
\%y, \%x}. All of this is computed modulo $2^{17}$, so a negative offset in
two's complement simply works.

\section{One Instruction, Cycle by Cycle}
\label{sec:step09-cycles}

Let us follow the instruction \texttt{addq 1, \%p, \%p} of \cc{01_hello.s}
(below) through the hardware. Here \reg{p} is \reg{1}, the instruction is at
address \texttt{0x08}, its code is \texttt{0x12110001} (opcode \texttt{12},
$z=1$, $y=1$, $\mathit{imm}=1$), and \reg{1} contains \texttt{0x20}. The values
in Table~\ref{tab:step09-addq} are taken from a simulation of this program
(Section~\ref{sec:step09-run}).

\begin{table}[ht]
\centering
\small
\begin{tabular}{@{}l >{\raggedright\arraybackslash}p{68mm} >{\raggedright\arraybackslash}p{44mm}@{}}
\toprule
state & during the cycle & at the rising edge at its end \\
\midrule
FETCH & RAM switch to instruction fetch; RAM address \cc{cu_ip} $=$ \texttt{0x08},
  \cc{RAM_FETCH} & SPRAM takes over the address \\ \addlinespace
DECODE & \cc{cu_ir} $=$ \texttt{0x12110001}; the decoder computes \cc{ALU_ADD},
  \cc{ALU_IMM}, \cc{s_reg} $=1$, \cc{b_reg} $=1$, \cc{a_imm} $=1$;
  \cc{decoder_en} $=1$ & \cc{instr_alu} takes over these values \\ \addlinespace
LOAD\_OPS & register file: \cc{addr_out1} $=$ \cc{b_reg} $=1$; \cc{alu.a} $=$
  \cc{a_imm} $=1$ & \cc{data_out1} $\gets \%1 =$ \texttt{0x20} \\ \addlinespace
EXECUTE & \cc{alu.op} $=$ \cc{ALU_ADD}, \cc{alu.b} $=$ \texttt{0x20},
  \cc{alu.s} $=$ \texttt{0x21}, ZF $=0$ & \cc{acc} $\gets$ \texttt{0x21} \\ \addlinespace
INCREMENT & \cc{alu.op} $=$ \cc{ALU_NOP}, so \cc{alu.s} $=$ \cc{acc} $=$
  \texttt{0x21}; register file: \cc{op} $=$ \cc{REG_WRITE}, \cc{addr_in} $=1$,
  \cc{data_in} $=$ \texttt{0x21}; \cc{cu_ip_next} $=$ \texttt{0x0c}
  & $\%1 \gets$ \texttt{0x21}, \cc{cu_ip} $\gets$ \texttt{0x0c}, next state
  FETCH \\
\bottomrule
\end{tabular}
\caption{\texttt{addq 1, \%p, \%p} in five clock cycles.}
\label{tab:step09-addq}
\end{table}

The load \texttt{movzbq (\%p), \%ch} before it, at address \texttt{0x04},
takes the same five states but uses other paths. In LOAD\_OPERANDS the
register file reads \reg{1} on port 0 (the address register). In EXECUTE the
RAM switch of the control unit has changed sides, and the data access sends
the address $\%1 + \mathit{offset} = $ \texttt{0x20} with \cc{RAM_FETCH} and
\cc{RAM_BYTE} to the RAM. In INCREMENT the RAM delivers the byte
\texttt{0x68} (an \texttt{h}), zero-extended to 64 bits by the byte logic of
Chapter~\ref{ch:step08}, and the register file writes it to \reg{2}. For
\texttt{movsbq} the decoder sets an additional signal \cc{instr_bus.sign_ext},
and the control unit replaces bits 63 to 8 of the byte by copies of its bit 7
before it goes to the register file.

At 12\,MHz and five cycles per instruction the ULM executes 2.4 million
instructions per second.

\begin{exercise}
\label{ex:step09-encode}
(a) The word \texttt{0x05FFFFFB} is at address \texttt{0x18} of
\cc{examples/01_hello}. Which instruction is it, and where does it jump? (b)
Encode \texttt{subq \%3, \%4, \%5}. Which register gets which difference?
\end{exercise}

\section{Running a Program}
\label{sec:step09-run}

This is the first example program:

\verbfile{examples/01_hello.s}

Line 5 puts a string of 65 bytes and a terminating zero byte into the data
segment: UTF-8 encodes the beer mug emoji in four bytes and each Chinese
character in three, and \texttt{\textbackslash n} is the newline character. The program sets
\reg{p} to the address of the string, loads one byte after another, and sends
it with \texttt{putc}, until it finds the terminating zero byte. The
comparison with zero is a subtraction whose result goes to \reg{0}: the result
itself is thrown away, but the zero flag remains. Assembled, the program is
this sequence of hex digits, \cc{examples/01_hello}:

\begin{terminal}
10100020202100001211000114020000040000033020000005FFFFFB0142000068656C6C6F
2C20776F726C642120F09F8DBA200AE4BDA0E5A5BDE4B896E7958C0A48616C6C6F2C205765
6C740A4772C3BCC39F652061757320556C6D203B2D290A00
\end{terminal}

(the file is one line; we broke it here). The first eight words are the
instructions; from address \texttt{0x20} on follows the string, byte by byte.
UTF-8 needs several bytes for the emoji and for the Chinese characters, and
the ULM does not care: it sends bytes, and the terminal puts them together.

\subsection*{On the board}

\begin{enumerate}
\item In \cc{09_ulm_on_ice}, \cc{make prog} writes the ULM to the board.
\item Start a serial terminal, for instance \cc{06_io_mirror/kermit-serial}
  (Chapter~\ref{ch:step06}; adjust the device in \cc{set line}). It talks at
  9600 baud.
\item Paste the contents of \cc{examples/01_hello} into the terminal. Each
  character comes back as an echo, since \cc{dev_io} copies every received
  byte to the output (Chapter~\ref{ch:step07}). The loader turns pairs of hex
  digits into bytes and ignores everything that is not a hex digit, such as
  line breaks. Meanwhile the display shows the byte being assembled.
\item Press Ctrl-D. This sends the byte 4; the loader sets \cc{done}, LED3
  goes on, the display goes dark, and the ULM runs. The text appears in the
  terminal, then the ULM halts: LED4 goes on, and the display shows the exit
  code \texttt{42} of \texttt{halt 0x42}.
\item To run the next program, press \cc{BTN1} and start again with step 3.
  The registers and the RAM are not cleared, the new program simply overwrites
  the old one from address 0 on.
\end{enumerate}

The other examples read numbers with \texttt{getc} and print results:
\cc{02_uint64_hex_io} works with hexadecimal numbers, \cc{03_uint64_dec_io}
with decimal ones, and \cc{04_factorial} computes factorials. Their source files
show how the ULM calls functions with \texttt{jmp \%y, \%x} and a stack in the
RAM.

\begin{pitfall}
A first version of \cc{test.sv} declared the reset as \cc{logic rst = BTN1;}.
In SystemVerilog this is a declaration with an \emph{initial value}, not an
\cc{assign}: the variable gets the value of \cc{BTN1} once, at time 0. yosys
turned it into a connection to \cc{BTN1}, so the button worked on the board;
Verilator follows the standard, so in the simulation \cc{rst} never changed.
(We checked both with a two-line module.) If a signal is meant to follow
another one, write \cc{logic rst;} and \cc{assign rst = BTN1;}, as
\cc{test.sv} does now. The testbench below relies on it: it presses
\cc{BTN1} before each program.
\end{pitfall}

\subsection*{In the simulation}

\cc{make sim} works as in all earlier steps. The testbench \cc{tb_test.cpp}
does not go through the serial line: at 9600 baud a byte takes $12\cdot10^6
\cdot 11 / 9600 = 13\,750$ clock cycles, far too many for a simulation.
Instead it plays the role of the UART receiver (Chapter~\ref{ch:step06})
itself, in the function that simulates one clock cycle:

\cpppart{firstline=39,lastline=70}{09_ulm_on_ice/tb_test.cpp}

Every fourth clock cycle it sets the outputs \cc{rx_ready} and \cc{rx_data}
of \cc{uart_rx0}, as if a byte had just arrived. Verilator makes this
possible: the C++ class of the design has members for the submodules, and the
testbench may overwrite their signals. In the other direction, the testbench
reads four signals of \cc{test.sv} that are marked with the comment
\cc{/* verilator public */}: \cc{cpu_halted}, \cc{cpu_exit_code}, and
\cc{cpu_putc_done} and \cc{cpu_putc_char}, which are \cc{putc_push_done} and
\cc{putc_char} of the CPU's I/O interface. Whenever \cc{cpu_putc_done} is 1,
\cc{dev_io} has just taken a character from a \texttt{putc}, and the
testbench appends it to \cc{output}. This is much faster than decoding the
serial signal on \cc{TX}.

The function \cc{run} presses \cc{BTN1} for a few cycles, sends a program as
hex digits followed by the byte 4, waits until the loader is done
(\cc{LED3}) and then until the CPU has halted, and compares the exit code and
the output with the expected ones. The testbench runs four programs this way:
\cc{examples/01_hello}, which it reads from the file, and three short
programs that check details of the control unit: that \texttt{putc 'A'} does
not write \reg{4}, that the backward jump \texttt{jnz} does not write
\reg{15}, and that \texttt{movsbq} extends the sign. Only the first program
goes into \cc{waveform.vcd}. The output, without the build steps (the lines
of \cc{01_hello} with Chinese characters and the emoji are left out here):

\outputfile{step09_make_sim.txt}

The loader is done after 802 clock cycles, and the ULM then sends the 65
bytes of the string and halts 2046 cycles later with exit code
\texttt{0x42}. Table~\ref{tab:step09-addq} was read off this simulation.

\begin{exercise}
\label{ex:step09-cycles}
Explain the 2046 cycles from the program and Figure~\ref{fig:step09-states}.
How long does it take the UART to send the 65 bytes at 9600 baud, with one
start bit, eight data bits and two stop bits each? Where do the bytes wait in
the meantime?
\end{exercise}

\section{Resources and Speed}
\label{sec:step09-resources}

\cc{make report} (Appendix~\ref{app:tools}) for the whole ULM:

\outputfile{step09_report.txt}

The ULM fills 57\,\% of the logic cells, all four SPRAMs, ten of the block
RAMs (eight for the registers, two for the FIFOs) and one DSP block. The
per-module part of the report is of little help this time: almost all modules
of this step have interfaces as ports, sv2v merges those into their parent,
and so the ALU, the register file and the control unit all disappear into
\cc{test}. The floorplan (\cc{make floorplan}) assigns logic cells to instances
by the names of the cells, and its table gives a better picture:

\outputfile{step09_floorplan.txt}

\textbf{The ALU} is the largest part. To find out why, we synthesized
\cc{dev_alu} alone, with a small wrapper module that connects its interface
to ports: 1375 LUTs, 128 carry cells, 67 flip-flops, one DSP. Without the two
shift operations it needs 540 LUTs. So the two 64-bit shifters cost about 835
LUTs: a shifter that can shift by any amount from 0 to 63 is a cascade of six
layers of 64 multiplexers (shift by 1, 2, 4, 8, 16, 32 or not), for each
direction. The 128 carry cells are two 64-bit carry chains, one for the
adder and one for the subtractor.

\textbf{The register file} gets 339 logic cells in the floorplan, for
something that lives in block RAMs. This is an artefact of how the floorplan
assigns cells: yosys names a cell after a net it is connected to, and many
LUTs of the ALU's input logic are named after the read data of the block RAMs.
Synthesized alone, with a small wrapper module, the register file gives

\outputfile{step09_rf_stat.txt}

that is, the eight block RAMs and three LUTs for the write enable
(\cc{op} is \cc{REG_WRITE} and \cc{addr_in} is not 0). Two details of
\cc{dev_reg_file.sv} keep it this small:
\begin{itemize}
\item \reg{0} reads as 0 without any logic. Address 0 is never written, and
  \cc{mem[0]} starts as 0 (the initial contents of a block RAM are part of the
  bitstream, Chapter~\ref{ch:step07}), so both copies of the memory always
  hold 0 there. A test like \cc{addr_out0 != 0 ? mem[...] : 0} would need a
  multiplexer behind the block RAM, one LUT for each of the 64 bits of each
  read port.
\item The attribute \cc{(* no_rw_check *)} on \cc{mem} tells yosys not to
  care about a read and a write of the same address at the same clock edge
  (Chapter~\ref{ch:step07}). Without it, yosys delays the write by one cycle
  (flip-flops for the 64 data bits, the address and the write enable) and adds
  a bypass, so that a read still sees the value the description promises. In
  the ULM, the register file is written only at the end of INCREMENT, and
  what its read ports deliver after that edge is not used: in the next
  instruction, the operands are read again in LOAD\_OPERANDS.
\end{itemize}
Exercise~\ref{ex:step09-rf} measures what each of them saves.

\textbf{The clock.} nextpnr reports a maximum frequency of 15.41\,MHz,
against the 12\,MHz of the board: the design works, with about 25\,\% of
margin. The critical path is this one (shortened):

\outputfile{step09_pnr_path.txt}

It starts at the flip-flops of the decoded instruction (\cc{instr_alu.a_imm}),
goes through the 64-bit adder of the ALU and the test whether the result is
zero (\cc{alu.zf}), and ends at the register in the decoder that decides
whether a conditional jump is taken (\cc{instr_cu.op}): 64.9\,ns of the
83.3\,ns that a clock cycle at 12\,MHz has. As in step 01, most of the time is
spent in the wires (45.0\,ns), not in the logic (19.9\,ns). Strictly speaking,
this path is never used within one clock cycle: the ALU computes only in
EXECUTE, and the decoder takes over its inputs only in DECODE, when the ALU
gets \cc{ALU_NOP} and \cc{alu.zf} comes from the register \cc{acc}. nextpnr
does not know the state machine, so it has to assume the worst case. Such a
path is called a \term{false path}.

\begin{exercise}
\label{ex:step09-rf}
In \cc{dev_reg_file.sv}, (1) remove \cc{(* no_rw_check *)}; (2) instead, keep
the attribute and read with \cc{reg_file.addr_out0 != 0 ? mem[reg_file.addr_out0]
: 0} (and the same for port 1). Predict what happens to the cell count of the
register file and of the whole design, then check with \cc{make report}. Does
the ULM still work (\cc{make sim})?
\end{exercise}

\begin{exercise}
\label{ex:step09-putc}
The testbench loads this program:
\begin{terminal}
10400007125400013141000002400000
\end{terminal}
It is \texttt{loadz 7, \%4}; \texttt{addq 1, \%4, \%5}; \texttt{putc 'A'};
\texttt{halt \%4}, and it halts with exit code \texttt{07}. Suppose
\cc{dev_cu.sv} wrote the register file in every INCREMENT state, that is,
without \cc{&& writes_reg}. Which exit code would the program give? Which
other instructions would change a register they do not name?
\end{exercise}

\section{What Could Be Improved}
\label{sec:step09-outlook}

The README of the repository calls the ULM on ICE a working prototype, written
while its author met SystemVerilog and FPGAs for the first time. It does what
it was built for: it runs the programs of the course. Having gone through all
nine steps, you are now in a position to improve it. Some ideas, from small to
large:

\begin{itemize}
\item \textbf{More instructions.} The instruction set has no \texttt{or},
  no \texttt{not}, no arithmetic shift right, no division, and only byte and
  quad word loads and stores. Each new instruction is an opcode in
  \cc{dev_decoder.sv}, perhaps an operation in \cc{dev_alu.sv}, and an entry
  in \cc{isa.txt}, so that the assembler knows it.
\item \textbf{Fewer logic cells.} The adder and the subtractor can share one
  carry chain, since $b - a = b + \neg a + 1$, as in the Mini-ALU of
  Chapter~\ref{ch:step05}. The shifters could shift by one bit
  per clock cycle, at the price of up to 63 cycles per shift. With more than
  40\,\% of the logic cells free, there is room for more instructions, such
  as a full 64-bit multiplication from four of the DSP blocks.
\item \textbf{More speed.} Five cycles per instruction is slow. The next
  instruction could be fetched while the current one is executed; this is
  \term{pipelining}, and every modern processor does it. It needs care: a
  jump changes which instruction comes next, and an instruction may need the
  result of the one before. And the clock could be faster: the iCE40UP5K has a
  PLL that can make a faster clock from the 12\,MHz, but first the critical
  paths must become shorter, and nextpnr must be told which paths are false
  ones, or the design must be changed so that they disappear.
\end{itemize}

We started with two buttons and an AND gate. We end with a computer that
we can describe completely, from the lookup table to the instruction set,
and that runs programs we write. Every step on the way was a circuit, and so
is the ULM: registers, multiplexers, adders and memory, connected by wires.

\solutions

\solution
(a) Opcode \texttt{05} is \texttt{jmp}. The 24-bit offset \texttt{0xFFFFFB}
is $-5$ in two's complement, so the target is $\texttt{0x18} - 5\cdot 4 =
\texttt{0x04}$, the instruction \texttt{fetch movzbq (\%p), \%ch}: this is the
\texttt{jmp fetch} that closes the loop. In hardware, the decoder appends two
zero bits, and \cc{dev_cu} adds the lowest 17 bits to \cc{cu_ip}: $\texttt{0x18}
+ (2^{17} - 20) \equiv \texttt{0x04} \pmod{2^{17}}$.

(b) The register form of \texttt{subq} has opcode \texttt{15} and format
\texttt{R\_R\_R} with $z = 5$, $y = 4$, $x = 3$: \texttt{0x15543000}. It
computes $\%5 \gets \%4 - \%3$ (the ALU computes \cc{b - a} with \cc{b} from
$y$ and \cc{a} from $x$).

\solution
The program executes \texttt{loadz} once, then for each of the 65 bytes the
six instructions \texttt{movzbq}, \texttt{addq}, \texttt{subq}, \texttt{jz}
(not taken), \texttt{putc}, \texttt{jmp}, then \texttt{movzbq}, \texttt{addq},
\texttt{subq} and the taken \texttt{jz} for the zero byte, and finally
\texttt{halt}. That is $1 + 6\cdot 65 + 4 + 1 = 396$ instructions of five
cycles, 1980 cycles. Each \texttt{putc} stays one cycle longer in EXECUTE:
\cc{dev_io} takes the character at the end of the first EXECUTE cycle and
answers with \cc{putc_push_done} in the second. That adds 65 cycles, and the
FETCH cycle in which the halt is noticed one more: 2046. That is 0.17\,ms at
12\,MHz. The UART needs 11 bit times per byte, so $65 \cdot 11 /
9600\,\mathrm{s} \approx 74.5$\,ms, more than 400 times as long. The bytes
wait in the output FIFO of \cc{dev_io}, which holds 512 bytes, and the ULM has
halted long before the terminal shows the text.

\solution
(1) Without \cc{(* no_rw_check *)}, the register file alone grows from 11 to
284 cells: 135 flip-flops and 141 LUTs for the delayed write and the bypass
come in addition to the block RAMs. The whole design grows from 3036 to 3228
logic cells, and its maximum frequency drops from 15.41 to 14.18\,MHz. (2)
With the test for address 0, the register file alone has 143 cells: the
multiplexers behind the block RAMs. The whole design needs 3070 logic cells
and reaches 14.98\,MHz. The saving in the whole design is smaller than in the
register file alone because nextpnr can put a LUT and an unrelated flip-flop
into the same logic cell. In both cases \cc{make sim} still passes all four
programs: the extra logic does what it should, it is just not needed in the
ULM.

\solution
Without \cc{writes_reg}, the defaults of the \cc{always_comb} block would
apply to \texttt{putc 'A'}: write \cc{alu.s} to register
\cc{instr_alu.s_reg}, which the decoder always cuts out of \cc{ir[23:20]}.
\texttt{putc 'A'} is \texttt{0x31410000}, so \cc{s_reg} is 4, and the ALU's
last result, $\mathtt{acc} = 7 + 1 = 8$, would be written to \reg{4}: exit
code \texttt{08}. (An earlier version of \cc{dev_cu.sv} had exactly this
mistake.) The same would happen for \texttt{halt imm} (register
$\lfloor\mathit{imm}/16\rfloor$) and for the relative jumps, taken or not: \texttt{jmp fetch} in \cc{01_hello} is
\texttt{0x05FFFFFB}, which would write \cc{acc} to \reg{15}, and so does every
backward jump, since its offset is negative and starts with the bits
\texttt{1111}. This is what the third program of the testbench checks. The
examples would survive it, because their \texttt{putc} constants are
\texttt{'\textbackslash n'} $=$ \texttt{0x0A}, which hits \reg{0}, and they do
not use \reg{15}.
